\par\medskip
\begin{center}\bfseries IV. Die Jahre 1799 und 1800. AgM. Scheda Ac. Allgemeine elliptische Funktionen.\end{center}
\begin{center}\textls[100]{a) Allgemeines. Das agM. nach dem 30. Mai 1799.}\end{center}
\indent »Da ich vor der Hand wohl nicht in die Ketten eines Amts treten werde
und in Braunschweig zu meinen Arbeiten zu wenig Hilfsmittel hatte, so fasste
ich den Entschluss, mich eine Zeit lang hierher nach Helmstedt zu begeben,
wo ich wohl bis zu Ostern bleiben werde«, so schreibt \textsc{Gauss} an seinen Freund
\textsc{Bolyai} aus Helmstedt am 16. Dezember 1799 und fährt dann fort: »Ich wohne
hier bei dem Professor \textsc{Pfaff}, den ich eben so sehr als einen trefflichen Geometer,
wie als guten Menschen und meinen warmen Freund verehre; ein
Mann von einem arglosen kindlichen Charakter, ohne alle die Leidenschaften,
die den Menschen so sehr entehren und bei Gelehrten so gewöhnlich sind.
— Da ich noch nicht einmal acht Tage hier bin, so kann ich noch nicht
entscheiden, wie ich übrigens hier zufrieden sein werde; der Ort selbst ist
affreux, die Gegenden umher werden gerühmt, Bequemlichkeiten des Lebens
muss man manche entbehren, der Ton unter den Studenten im Ganzen soll
ziemlich roh sein, unter den Professoren, die ich habe kennen lernen, sind
artige Männer«\footnote{\textit{Briefwechsel Gauss-Bolyai}, S. 36.}. Dass \textsc{Gauss} ohne die »Ketten eines Amts« gerade in der
für die wissenschaftliche Arbeit fruchtbarsten Lebensperiode, bis zur Vollendung
des dreissigsten Jahres, sich ganz seinen Untersuchungen hingeben konnte,
verdankte er dem feinen Verständnis, das Herzog \textsc{Carl Wilhelm Ferdinand}\par
\clearpage\SourceLeaf{148}\ArithmeticHead{63}{R}{IVa. DAS AGM. NACH DEM 30. MAI 1799.}
\noindent für seine geistige Grösse bekundete. \textsc{Gauss} selbst hat dem väterlichen Freunde
in seiner Widmung der \textit{Disqu. arithm.} (Werke I, S. 3, 4) seine Dankbarkeit
bezeugt, aber auch von der Nachwelt gebührt dem unglücklichen Fürsten\footnote{Er starb am 10. November 1806 auf der Flucht vor dem Feinde an den Folgen einer Verwundung, die er als Führer der preussischen Armee in der unglücklichen Schlacht von Auerstädt erlitten hatte.}
ein Zoll dankbarer Anerkennung dafür, dass er dem grossen Sohne Braunschweigs
die freie, von keiner Sorge um des Lebens Notdurft behinderte Entfaltung
seiner Geistesgaben ermöglicht hat.

Über die Beschäftigung von \textsc{Gauss} in den Sommermonaten des Jahres 1799
besitzen wir keinerlei Aufzeichnungen, erst der September bringt eine Tagebuchnotiz
(99.), die von bedeutenden Fortschritten in den Prinzipien der Geometrie
berichtet (vergl. \textsc{Stäckel}, Abh. 4 dieses Bandes S. 23), und im November
meldet die Tagebuchnotiz 100. noch aus Braunschweig, dass er ȟber das
arithmetisch-geometrische Mittel vieles neue entdeckt« habe. Es ist wohl anzunehmen,
dass diese Notiz nur besagen will, dass \textsc{Gauss} vor der Abreise
nach Helmstedt (siehe oben) die in den vorhergehenden Monaten gemachten
Entdeckungen zusammengefasst und in die eben damals begonnene Scheda Ac
einzutragen begonnen hat. Die Leisteaufzeichnungen (Werke X, 1, S. 177—180)
dürften wohl etwas früher, noch während der eigentlichen Arbeit geschrieben
sein, denn es sind ziemlich zusammenhanglose Formeln, ohne jeden
Text.

Der ausschlaggebende Gedanke, der \textsc{Gauss} durch die Beobachtung des
30. Mai 1799 suggeriert worden ist, war, statt des agM. selbst, seinen
reziproken Wert zu betrachten, denn der reziproke Wert von \(M(1,\sqrt2)\) wird,
abgesehen von dem Faktor \(1/\pi\), durch die lemniskatische Periode \(\varpi\) gegeben.
In der 1800 begonnenen Abhandlung über das agM. (Werke III, S. 366,
367)\footnote{Diese Abhandlung sowie alles aus dem Nachlass, was auf das agM. und die allgemeine Theorie der elliptischen Funktionen Bezug hat, hat \textsc{Harald Geppert} in Nr. 225 von \textsc{Ostwalds} Klassikern übersetzt und durch verbindenden Text mit ergänzenden Einschaltungen zu einem trefflich lesbaren Ganzen gestaltet.} kündigt \textsc{Gauss} diesen Gedanken mit den eine innere Bewegung verratenden
Worten an: »Da ihre Koeffizienten [nämlich die der Reihe für
\(M(1,1+x)\), siehe oben S. 42] aber kein durchsichtiges Gesetz erkennen lassen,
gehen wir über diese Reihe hinweg und schlagen einen andern Weg ein, der
zu einem glücklicheren Erfolge führen wird.«\par
\clearpage\SourceLeaf{149}\ArithmeticHead{64}{L}{L. SCHLESINGER, ÜBER GAUSS’ FUNKTIONENTHEORETISCHE ARBEITEN.}
\indent In den Leistenotizen finden wir die Reihe, die den reziproken Wert von
\(M(1,p)\) darstellt und nach Potenzen von \(\nu=\sqrt{1-p^2}\) fortschreitet (Werke X, 1,
S. 177, vergl. auch die Bemerkungen, ebenda S. 261—266),
\begingroup\makeatletter\tagsleft@true\let\veqno\leqno\makeatother
\begin{equation*}
\frac1{M(1,p)}=1+\left(\frac12\right)^2\nu^2+\left(\frac{1\,.\,3}{2\,.\,4}\right)^2\nu^4+\cdots,\tag*{(1)}
\end{equation*}
\noindent und gleich daneben die Entwicklung für den Quadranten
\begin{equation*}
q=\int_0^1\frac{(1-\nu^2x^2)\,dx}{\sqrt{1-x^2\,.\,1-\nu^2x^2}}
=\frac\pi2\left(1-\left(\frac12\right)^2\frac{\nu^2}{1}-\left(\frac{1\,.\,3}{2\,.\,4}\right)^2\frac{\nu^4}{3}-\cdots\right)\tag*{(2)}
\end{equation*}
\noindent der Ellipse mit den Halbachsen \(1,p\), für \(p<1\). Durch Reihenvergleichung
ergeben sich dann die Beziehungen
\begin{equation*}
q=\frac\pi2\frac1{M(1,\sqrt{1-\nu^2})}+\nu\frac{dq}{d\nu}\hbox{\footnotemark},\tag*{(3)}
\end{equation*}
\footnotetext{Auf einem Blatte (Ff) finden sich die folgenden Formeln:
\[
\begin{aligned}
\text{»}E&=1-\frac1{2\,[2]}xx-\frac12\,.\,\frac14\,.\,\frac12\,.\,\frac34x^4-\frac{1\,.\,1\,.\,3\,.\,3\,.\,5}{2\,.\,2\,.\,4\,.\,4\,.\,6\,.\,6}x^6\ldots\\
F&=1+\frac1{2\,.\,2}xx+\frac12\,.\,\frac12\,.\,\frac34\,.\,\frac34x^4+\frac{1\,.\,1\,.\,3\,.\,3\,.\,5\,.\,5}{2\,.\,2\,.\,4\,.\,4\,.\,6\,.\,6}x^6\ldots\text{«}
\end{aligned}
\]
Es wird dann durch Differentiation von \(F\) und Reihenvergleichung die Formel
\[
(1-xx)\left(F+\frac{x\,\partial F}{\partial x}\right)=E
\]
abgeleitet, die mit der Formel (3) unseres Textes gleichwertig ist. Auch steht auf derselben Seite dieses Blattes die Formel
\[
F(x)=\frac1{M(1+x,\,1-x)}\mathpunct{\vcenter{\hbox{.}}}
\]
}
\begin{equation*}
-\frac2\pi\frac{dq}{d\nu}=\frac\nu{M(1,\sqrt{1-\nu^2})}+(\nu^2-1)\frac d{d\nu}\frac1{M(1,\sqrt{1-\nu^2})}\mathpunct{\vcenter{\hbox{,}}}\tag*{(4)}
\end{equation*}
\endgroup
\noindent die historisch insofern von besonderem Interesse sind, als sie Beispiele für
die »Relationen zwischen verwandten Funktionen« darstellen, für die durch
die angegebenen Potenzreihen gelieferten speziellen Fälle der Reihe \(F(\alpha,\beta,\gamma,x)\).

Die Reihenentwicklungen für \(q\) und \(\frac{dq}{d\nu}\) dürfte \textsc{Gauss} der \textsc{Eulerschen}
Abhandlung \textit{Animadversiones in rectificationem ellipsis} (Opuscula varii argumenti
II, 1750, S. 121 ff.; Opera, ser. I, vol. 20, S. 21) entnommen haben, wo diese
Entwicklungen in derselben Form auftreten wie in den Leistenotizen.\par
\clearpage\SourceLeaf{150}\ArithmeticHead{65}{R}{IVa. DAS AGM. NACH DEM 30. MAI 1799.}
\HistoricalFoot{X\,{\scriptsize 2} Abh. 2.}{9}
\indent Im höchsten Grade überraschend wirkt es aber, wenn nun \textsc{Gauss} völlig
unvermittelt durch die Gleichung
\begingroup\makeatletter\tagsleft@true\let\veqno\leqno\makeatother
\begin{equation*}
\frac\nu{M(1,\sqrt{1-\nu^2})}=\left(2z^{\frac12}+2z^{\frac92}+\cdots\right)^2\tag*{(5)}
\end{equation*}
\noindent den Zusammenhang zwischen dem agM. und den Reihen herstellt, deren Exponenten
die Quadratzahlen sind.

Wie er die Beziehung solcher Reihen zum agM., insbesondere wie er
jene Gleichung (5) gefunden haben mag, lässt sich nur vermuten. Es ist anzunehmen,
dass \textsc{Gauss} bemerkt hat, dass die Quadrate der Reihen\footnote{Von den drei Reihen \(p,q,r\) kommt in diesen Leistenotizen nur die letzte \(r\) explizite vor; \(p\) erscheint nur in den Differentiationsformeln (8), (9) Werke X, 1, S. 177, aber gerade das zeigt, dass die drei Reihen \textsc{Gauss} durchaus geläufig waren, und dass wir es in den in Rede stehenden Notizen mit flüchtigen, bei der Arbeit gemachten Aufzeichnungen zu tun haben.}
\begin{equation*}
\left\{\begin{aligned}
p(x)&=1+2x+2x^4+2x^9+\cdots,\\
q(x)&=1-2x+2x^4-2x^9+\cdots,\\
r(x)&=2x^{\frac14}+2x^{\frac94}+2x^{\frac{25}4}+\cdots,
\end{aligned}\right.\tag*{(\(\ast\))}
\end{equation*}
\noindent den Gesetzen des agM. gehorchen, wenn man \(x\) durch \(x^2,x^4,x^8,\ldots\) ersetzt
(vergl. weiter unten die Gln. (28), (29)), und dass er daraufhin die Gleichung
(5) und die analogen einfach angesetzt hat, indem er sich die Grösse \(z=\sqrt x\)
direkt durch jene Gleichung (5) definiert denkt. Für die letztere Annahme
sprechen die von \textsc{Gauss} a. a. O. aufgezeichneten Gleichungen
\begin{equation*}
\left\{\begin{aligned}
\nu+\frac14\nu^3+\frac9{64}\nu^5+\cdots&=4z+8z^5+\cdots\\
\nu&=4z-16z^3+56z^5-\cdots,
\end{aligned}\right.\tag*{(6)}
\end{equation*}
\noindent denen sich noch die Darstellung des Ellipsenquadranten \(q\)
\begin{equation*}
\frac2\pi q=1-4z^2+20z^4+663z^6+\cdots\tag*{(7)}
\end{equation*}
\endgroup
\noindent zugesellt\footnote{Alle diese Formeln enthalten zahlreiche Schreibfehler, die wohl daher rühren, dass \textsc{Gauss}, wie auch sonst häufig, gewisse Gleichungen einfach aus dem Gedächtnis hingeschrieben haben dürfte.

Wir bemerken noch, dass die drei Reihen \(p,q,r\), wie sie in den Gleichungen (\(\ast\)) definiert sind, in der Variablen \(z=x^{\frac12}\) geschrieben in der Scheda Af (1801) auftreten, siehe Werke X, 1, S. 210, Art. 3.}.

Für die Herleitung aller dieser Formeln, insbesondere der Gleichungen
(6), (7), kommt die ausserordentliche Gewandtheit in Betracht, die sich \textsc{Gauss}\par
\clearpage\SourceLeaf{151}\ArithmeticHead{66}{L}{L. SCHLESINGER, ÜBER GAUSS’ FUNKTIONENTHEORETISCHE ARBEITEN.}
\noindent schon frühzeitig in dem Operieren mit Reihen, deren Inversion (durch die
\textsc{Lagrangesche} Formel, vergl. oben) und Transformation angeeignet hat. Dass
er aber die Bedeutung der Grösse \(z\) gleich erkannt hat, indem er mit ihrer
Hilfe in der Gleichung (7) den Ellipsenquadranten als Funktion des Moduls \(\nu\)
uniformisiert, ist der vorhandenen Kraft des mathematischen Genius zuzuschreiben,
dessen geheimnisvolles Wirken wohl nirgends in so überraschender
Weise zu verspüren ist, wie eben in diesen Jugendarbeiten von \textsc{Gauss}.

Als der wesentlichste Fortschritt, der in den geschilderten Leisteaufzeichnungen
enthalten ist, dürfte \textsc{Gauss} der Zusammenhang erschienen sein, der
das agM. mit einer »klassischen Grösse«, nämlich dem Ellipsenquadranten, verknüpft.
Den Beweis für die Beobachtung des 30. Mai kann er nunmehr
den Formeln entnehmen, die in der Scheda Ac auf S. 7 zwischen geometrischen
Entwicklungen\footnote{Diese Scheda hat auf dem Titelblatt die Aufschrift »Varia Novbr. 1799«. Dann folgt »Imprimis de integrali \(\displaystyle\int\frac{d\varphi}{\sqrt{1+\mu\mu\sin\varphi^2}}\)«, was aber, wie auch Schrift und Tinte zeigt, nicht gleichzeitig mit der Aufschrift selbst, also offenbar später angebracht worden ist. Die Seiten 1—8 enthalten »Grundlehre oder Analysis des Raumes« (räumliche Koordinaten und sphärische Trigonometrie).} aufgezeichnet sind. Dort finden wir nämlich in
sorgfältiger Schrift und von Strichen umrahmt die Reihen (abgedruckt Werke
X, 1, S. 184)
\begingroup\small\makeatletter\tagsleft@true\let\veqno\leqno\makeatother
\begin{equation*}
\frac\varpi\pi=\left(1+\left(\frac12\right)^2\frac12+\left(\frac12\,.\,\frac34\right)^2\,.\,\frac14+\left(\frac12\,.\,\frac34\,.\,\frac56\right)^2\frac18+\cdots\right)\sqrt{\frac12}\tag*{(\(\alpha\))}
\end{equation*}
\begin{equation*}
{}=\left(1+\frac{1\,.\,2\,.\,3}{2\,.\,4\,.\,4}\,.\,\frac19+\frac{1\,.\,2\,.\,3\,.\,5\,.\,6\,.\,7}{2\,.\,4\,.\,4\,.\,6\,.\,8\,.\,8}\,.\,\frac1{81}+\frac{1\,.\,2\,.\,3\,.\,5\,.\,6\,.\,7\,.\,9\,.\,10\,.\,11}{2\,.\,4\,.\,4\,.\,6\,.\,8\,.\,8\,.\,10\,.\,12\,.\,12}\,.\,\frac1{729}\ldots\right)\sqrt{\frac23}\tag*{(\(\beta\))}
\end{equation*}
\begin{equation*}
{}=\left(1+\frac14\,.\,\frac34\,.\,\frac19+\frac14\,.\,\frac34\,.\,\frac58\,.\,\frac78\,.\,\frac1{81}+\frac{1\,.\,3\,.\,5\,.\,7\,.\,9\,.\,11}{4\,.\,4\,.\,8\,.\,8\,.\,12\,.\,12}\,.\,\frac1{729}\ldots\right)\sqrt{\frac23}\mathpunct{\vcenter{\hbox{.}}}\tag*{(\(\gamma\))}
\end{equation*}
\endgroup
\noindent Die Reihe (\(\alpha\)) gibt, wenn man sie mit der Entwicklung (1) für \(\frac1{M(1,p)}\) vergleicht,
unmittelbar den Beweis der Gleichung des 30. Mai, und das war
offenbar der Grund, weshalb \textsc{Gauss} diese ihm schon seit 1797—98 bekannten
Entwicklungen hier noch einmal mit besonderer Sorgfalt eingetragen hat\footnote{Die Reihe (\(\alpha\)) findet sich auch (mit der älteren Bezeichnung \(\Pi\) statt \(\frac12\varpi\) für die lemniskatische Periode) auf der Rückseite des Titelblatts von \textit{Leiste} (siehe Werke X, 1, S. 169, Art. 5). Ebenda ist auch die Reihe (\(\beta\)) angedeutet und mit ihrer Hilfe der Wert von \(\Pi\) auf 15 Dezimalstellen berechnet. Dieselbe Reihe (\(\beta\)) ist auch im \textit{Tagebuch} (Notiz 91 a) mit der Datierung Juli [1798] aufgezeichnet. Die Reihe (\(\alpha\)) tritt auch bei S. 54 des \textit{Leiste} (Werke X, 1, S. 149) als absolutes Glied der Entwicklung von \((1+\sin^2\varphi)^{-\frac12}\) nach den Kosinus der Vielfachen von \(\varphi\) auf.}.
